\chapter{Subquadratic Sparse Attention}
\label{ch:ssa}

\newthought{Subquadratic Sparse Attention (SSA)} is the mechanism behind
the SubQ company. This chapter explains, with the limited information
that has been made public as of mid-2026, what it is, why it lives in the
bottom-right quadrant of Chapter~\ref{ch:quadrant}, and what is honestly
known versus what is still proprietary.\sidenote{Primary source for this
chapter: Subquadratic's technical post \citep{subq2026ssa}, supplemented
by the company's blog and demos. Where claims are vendor-reported, that
is said out loud.}

\section{The mechanism, in one sentence}

For each query, score it against every key in the past, but only
\emph{actually compute} a small data-dependent subset of those scores —
specifically, the ones likely to be largest — and treat the rest as
zero.

That subset is chosen by a learned selection function that does not require
materialising the full \(n \times n\) attention matrix to know which
entries matter. The selection function's cost is sublinear or linear in
\(n\), not quadratic.

\section{Why this isn't trivial}

The naive way to ``only score the top-\(k\) keys'' is to score \emph{all}
keys, sort, and keep the top \(k\). That is \(\bigO(n)\) per query and
\(\bigO(n^2)\) per layer — exactly what we are trying to avoid.

The non-trivial bit is finding the top \(k\) without doing the all-pairs
dot product first. This is the same problem as approximate nearest
neighbour search, and like ANN it is solved by some combination of:
\begin{itemize}
\item \textbf{Indexing the keys.} Build a structure (LSH, IVF, HNSW,
learned routing tree) that, given a query, returns a small candidate set
of keys in $\bigO(\log n)$ or $\bigO(\sqrt n)$ work. Then score only
those candidates.
\item \textbf{Hierarchical selection.} Group keys into clusters at training
time; route a query to one or a few clusters; do dense attention within
the chosen clusters.
\item \textbf{Learned routing.} A small network looks at the query and
predicts which keys are likely to be relevant, without computing the
actual dot products.
\end{itemize}
\citet{subq2026ssa} describe their selection as content-dependent and
hardware-friendly, but the specific algorithm is proprietary.

\section{What this buys you, on paper}

\begin{itemize}
\item \textbf{Cost.} If the selection function returns $\bigO(\log n)$ or
$\bigO(1)$ keys per query and itself costs $\bigO(\log n)$ per query, the
per-token attention cost is essentially constant in $n$. Total per-layer
cost becomes $\bigO(n \log n)$ rather than $\bigO(n^2)$.
\item \textbf{Routing fidelity.} Unlike sliding-window or state-space
models, the selected keys can be \emph{anywhere in the past}. A relevant
fact 800K tokens ago can be retrieved as easily as one 8 tokens ago, as
long as the selection function thinks to look at it.
\item \textbf{Drop-in.} The mechanism is meant to slot into a standard
transformer's attention block, allowing existing transformer training and
inference infrastructure to be reused.
\end{itemize}

\section{What it costs in correctness}

The honest version is that any sparse-attention mechanism trades off
\emph{recall}: the selection function will sometimes miss a key that
\emph{should} have been attended to, and the model will not see that
information. Two failure modes are worth naming.

\paragraph{Selection-function false negatives.} If the selector thinks key
\(s\) is irrelevant when it actually is, the dependent token \(t\) gets
no signal from \(s\). The model's behaviour will quietly degrade on tasks
where the relevant key looks superficially unrelated to the query.

\paragraph{Training-distribution dependence.} The selection function is
learned from data. If a long-context use case is sufficiently
out-of-distribution, the selector may pick the wrong keys. This is one
of the main reasons SubQ-style architectures need extensive
long-context-specific training, not just a long-context evaluation tail
on top of a standard transformer.

\section{Where SSA sits relative to the alternatives}

In the 2$\times$2 frame from Chapter~\ref{ch:quadrant}: bottom-right.
Linear-cost (or near-linear-cost), content-dependent routing. Distinct
from:

\begin{itemize}
\item Fixed-pattern sparse (Chapter~\ref{ch:sparse}) — same asymptote,
position-fixed routing.
\item State-space models (Chapter~\ref{ch:ssm}) — same asymptote,
content-dependent dynamics but \emph{lossy compression} of the past into
a finite state.
\item Linear attention / Hyena (Chapter~\ref{ch:linattn}) — better-than-
quadratic asymptote, but expressivity-limited routing kernels.
\item Hybrids (Chapter~\ref{ch:hybrids}) — mostly position-fixed mixers
plus a few full-attention layers; complementary to SSA, not competitive.
\end{itemize}

\keyidea{SSA's claim is that it occupies a quadrant nobody else can
practically reach: linear-cost \emph{and} content-dependent. If the
claim holds at frontier scale, it is a genuine architectural advance,
not an implementation trick.}
